(a) By Theorem \claimref{transport.existence-classification-type-c} with Lemmas \claimref{transport.centre} and \claimref{transport.sign}, an odd-dimensional algebra is type C exactly when $p' - q' \equiv 1 \pmod 4$ and complex type exactly when $p' - q' \equiv 3 \pmod 4$; these exhaust the odd residues. For $A_s$, $p' - q' = p - q - 1$, which is $\equiv 1$ iff $p - q \equiv 2$ and $\equiv 3$ iff $p - q \equiv 0 \pmod 4$. For $A_t$, $p' - q' = p - q + 1$, which is $\equiv 1$ iff $p - q \equiv 0$ and $\equiv 3$ iff $p - q \equiv 2$. Since $n$ is even, $p - q$ is even, so exactly one case obtains for each extension and the two extensions fall on opposite sides. Simplicity of the complex-type algebra: its centre $\C$ is a field, and a semisimple algebra with a field for centre has one Wedderburn block.

(b) The corner isomorphism $\varepsilon_+ A' \cong \Cl^0(A')$ is Lemma \claimref{transport.corner-even-subalgebra}, and $\varepsilon_- A' \cong \varepsilon_+ A'$ through the grade involution, which exchanges the ideals (Proposition \claimref{transport.ascent-renders-grade-involution-inner}). Lemma \claimref{transport.even-subalgebra-one-dimensional-extension} identifies $\Cl^0(p,q{+}1) \cong \Cl(p,q)$ and $\Cl^0(p{+}1,q) \cong \Cl(q,p)$. In the timelike case, $p - q \equiv 0 \pmod 4$ gives $q - p \equiv -(p-q) \pmod 8$, and the classification (Theorem \claimref{transport.classification-real-clifford-algebras}) assigns $\Cl(q,p)$ and $\Cl(p,q)$ the same type and dimension when $p - q \equiv 0 \pmod 4$ (residues $0$ and $4$ are each fixed by negation mod 8), so $\Cl(q,p) \cong \Cl(p,q)$ there (the script exercises the timelike form of Lemma \claimref{transport.even-subalgebra-one-dimensional-extension} on $(1,3)$ and $(2,0)$).

(c) The spacelike case: the repair generator $\nu$ has $\nu^2 = +1$ and the extension generator $e_{n+1}$ has square $-1$, so the Gram matrix of the pair is $\mathrm{diag}(+1,-1)$, of signature $(1,1)$: a hyperbolic plane. Then $f = \tfrac12(\nu + e_{n+1})$ satisfies $f^2 = \tfrac14(\nu^2 + e_{n+1}^2 + \nu e_{n+1} + e_{n+1}\nu) = 0$, likewise $f^* = \tfrac12(\nu - e_{n+1})$, and $ff^* + f^*f = \tfrac14\bigl((\nu + e_{n+1})(\nu - e_{n+1}) + (\nu - e_{n+1})(\nu + e_{n+1})\bigr) = \tfrac12(\nu^2 - e_{n+1}^2) = 1$. The quadratic space is $\R^{p,q} \perp \langle e_{n+1}\rangle \perp \langle \nu \rangle = \R^{p+1,q+1}$ whichever order the two directions are adjoined in, so the algebra is $\Cl(p{+}1,q{+}1)$ on the nose, and Theorem \claimref{transport.rule-1-morita-trivial} reads it as $\Mat_2(A)$ through the ladder pair, while Lemma \claimref{transport.matrix-units-peirce-isomorphism} and Theorem \claimref{transport.rule-2-ideal-fusion-sharp-form} read it as $\Mat_2(\Cl^0(A_s))$ through $\varepsilon_\pm$ and $\nu$; by (b) the two coefficient algebras are the same algebra. The residue of $\Cl(p{+}1,q{+}1)$ is $p - q$. The timelike case is the same computation with the signs of $\nu^2$ and $e_{n+1}^2$ exchanged and $f = \tfrac12(e_{n+1} + \nu)$, $f^* = \tfrac12(e_{n+1} - \nu)$.

(d) Theorem \claimref{transport.rule-2-ideal-fusion-sharp-form}(3) gives the isomorphism of the two repairs and the residues $d' \pm 1$ of the fused algebras with $d' = d \mp 1$ the residue of the extension; the same-sign repair has residue $d \mp 2$. The adjoined plane has Gram matrix $\pm\,\mathrm{diag}(1,1)$, definite, so $(a e_{n+1} + b\nu)^2 = \pm(a^2 + b^2) \neq 0$ for $(a,b) \neq 0$: the adjoined plane carries no null vector, hence no ladder pair, so the adjunction is not a $(1,1)$ ascent; the algebra it produces is nevertheless $\Mat_2(A)$, and does contain ladder pairs, just none inside the adjoined plane.

The script verifies (a) on every even-dimensional $(p,q)$ with $2 \leq n \leq 6$ by computing the volume element's centrality and square directly, and verifies (c) and (d) on the same grid: the obstructed extension is type C, the opposite-sign repair is at residue $d$ with a hyperbolic adjoined plane carrying a ladder pair, and the same-sign repair is two residues away with a definite plane.
